\documentclass[11pt,a4paper]{article}

\usepackage[margin=1in]{geometry}
\usepackage{amsmath,amssymb}
\usepackage{booktabs}
\usepackage{enumitem}
\usepackage{microtype}

\setlist[enumerate]{itemsep=0.35em,topsep=0.4em}
\setlist[itemize]{itemsep=0.25em,topsep=0.35em}

\title{Project plan:\\hydrodynamic entropy of a learned\\two-dimensional flow}
\author{}
\date{October 2026}

\begin{document}
\maketitle

\begin{abstract}
The quantity to compute is the hydrodynamic entropy of Verma and Chatterjee
(Phys.\ Rev.\ Fluids \textbf{7}, 114608, 2022) and of Verma, Stepanov, and Delache
(Phys.\ Rev.\ E \textbf{110}, 055106, 2024): the Shannon entropy of the Fourier
modal energies of a velocity field.
This plan trains a small network to advance a two-dimensional flow by one time
step, then asks whether the learned flow reproduces the entropy trajectories
those papers measured in the solver.
The entropy is never computed from weights, layers, or receptive fields.
\end{abstract}

\section{What the papers actually compute}

Both papers use the same definition.
For a velocity field with Fourier coefficient \(\mathbf{u}(\mathbf{k})\), the modal energy and its probability are
\begin{equation}
  E(\mathbf{k}) = \tfrac12 |\mathbf{u}(\mathbf{k})|^2,
  \qquad
  p_{\mathbf{k}} = \frac{E(\mathbf{k})}{\sum_{\mathbf{k}} E(\mathbf{k})},
  \qquad
  S_H = -\sum_{\mathbf{k}} p_{\mathbf{k}}\log_2 p_{\mathbf{k}}.
  \label{eq:SH}
\end{equation}
The sum runs over the active, dealiased modes.
If the energy is spread equally over \(M\) modes, \(S_H = \log_2 M\).
If one mode holds all the energy, \(S_H = 0\).
The formula uses amplitudes only. Phases do not enter.

For an isotropic spectrum the 2024 paper replaces the modal sum by a shell sum,
\(E(k) = \sum_{k-1 < |\mathbf{k}'| \le k} E(\mathbf{k}')\), with the same formula for \(S_H\).
The 2022 paper keeps the modal sum, because the structures it studies (a shear layer, a vortex dipole) are not isotropic.
This project uses the modal sum as the primary number, and the shell sum as a second number when the flow is forced and roughly isotropic.

\subsection{2022, two-dimensional Euler}

The system is incompressible Euler flow: no viscosity and no force.
Thermodynamic entropy is constant, because there is no heat production.
Hydrodynamic entropy is not.
On a \(512^2\) pseudospectral grid, three initial conditions reach three different large-scale states, and in each case \(S_H(t)\) falls and then levels off:

\begin{center}
\begin{tabular}{llrr}
\toprule
Run & Asymptotic structure & \(S_H\) & \(\log_2 M\) \\
\midrule
A & vortex--antivortex pair & 4.9 & \(\approx 16.5\) \\
B & unidirectional shear & 1.2 & \(\approx 16.5\) \\
C & four vortex pairs & 3.1 & \(\approx 16.5\) \\
\bottomrule
\end{tabular}
\end{center}

Run B is the cleanest illustration. The modes \(\mathbf{u}(0,\pm 1)\) hold about 99\% of the energy, so \(S_H \approx 1\).
The fall in \(S_H\) is the inverse cascade: energy moves toward a few small-\(k\) modes.
Those modes stay out of equilibrium. Their transfer \(T(k)\) and flux
\(\Pi(k) = -\int_0^k T(k')\,dk'\) remain nonzero, and \(\Pi(k)\) is negative at small \(k\).
At large \(k\) the spectrum settles near Kraichnan's equilibrium form
\(E(k) = k/(\beta + \gamma k^2)\), which is enstrophy equipartition, \(E(k)\sim k^{-1}\).
Three-dimensional Euler does the opposite: \(S_H\) rises toward \(\log_2 M\) and the flux vanishes.
The initial condition chooses the asymptotic value. A \(\delta\)-correlated start thermalizes and \(S_H\) stays at the maximum.

\subsection{2024, what that number means}

The follow-up paper uses the same \(S_H\) on forced Navier--Stokes turbulence, on the Ising model, and on the time-dependent Ginzburg--Landau equation, and states the properties that matter here.

\begin{itemize}
  \item \(S_H\) measures how the macroscopic energy is spread across scales. It is not thermodynamic entropy, and the two are not added.
  \item It is nonextensive. The ceiling is \(\log_2 M\), and a turbulent spectrum does not even reach that ceiling.
  \item For a Kolmogorov spectrum \(E(k)\propto k^{-5/3}\), logarithmically binned shells give \(S_H\) between about 3 and 4, almost independent of grid size. Their \(512^3\) forced run saturates near \(3.8\).
  \item A signal whose \(S_H\) is far from that range is not that kind of turbulence. Equipartition sits at \(\log_2 M\). A field that has collapsed onto one mode sits at \(0\), as in the Ginzburg--Landau run that relaxes to \(\phi = 1\).
  \item Because phases are discarded, \(S_H\) does not determine the energy transfer. The flux has to be computed separately.
\end{itemize}

\section{Where the network enters}

The papers already produce \(S_H(t)\) by integrating the fluid equations.
The project inserts a learned map in place of one time step of that integrator,
\[
  \omega_{t+\Delta t} = \mathcal{N}_\theta(\omega_t),
\]
and asks whether the entropy trajectory of the learned dynamical system is the trajectory in the papers.

The field that enters Eq.~\eqref{eq:SH} is the velocity (or the vorticity, converted to velocity through the stream function) of the flow, at a snapshot produced either by the solver or by rolling \(\mathcal{N}_\theta\) forward.
A matched one-step pixel error is not the result.
The result is whether \(S_H(t)\) under the learned step follows the paper: down to a small, initial-condition-dependent plateau for 2D Euler, or flat and of order a few bits for forced two-dimensional turbulence.

This is the same mechanism the papers identify.
In 2D Euler, order appears because nonlinear transfer piles energy into a few large-scale modes, and \(S_H\) falls.
A network that blurs the field also lowers \(S_H\), but it does so by destroying small scales, which is the Ginzburg--Landau collapse toward \(S_H = 0\), not the Euler condensate.
A network that thermalizes raises \(S_H\) toward \(\log_2 M\), which is the three-dimensional Euler outcome and the wrong attractor in 2D.
The flux sign separates these three outcomes. Only the Euler condensate has negative \(\Pi(k)\) at small \(k\) while \(S_H\) falls to a value of order 1, far below \(\log_2 M\) and above 0.

\section{Question and predictions}

\textbf{Question.}
Does a one-step neural emulator of two-dimensional flow reproduce the hydrodynamic-entropy trajectories of the 2022 and 2024 papers, including their dependence on the initial condition and on whether the flow is inviscid or forced?

Predictions, each tied to a result in the papers:

\begin{enumerate}
  \item \textbf{Euler, emergence of order.}
  On the three families of initial data from the 2022 paper, solver and emulator both have a falling \(S_H(t)\) that saturates far below \(\log_2 M\).
  The shear-like initial condition saturates lowest, as Run B does.

  \item \textbf{The plateau remembers the initial condition.}
  Training on one family and rolling out another does not produce that family's plateau.
  The 2022 states are not a unique equilibrium.

  \item \textbf{Forced turbulence, a steady entropy.}
  With steady forcing and viscosity, \(S_H(t)\) of a long rollout fluctuates about a constant of a few bits and does not grow when the grid is refined from \(64^2\) to \(128^2\).
  That is the nonextensivity result of the 2024 paper, in two dimensions.

  \item \textbf{Pixel loss alone leaves the trajectory.}
  One-step mean-squared error can look small while the rollout entropy goes to \(0\) (over-smoothed) or climbs toward \(\log_2 M\) (scrambled phases, equipartition).
  Matching the shell energies, and matching \(\Pi(k)\) at low \(k\), is what holds the rollout on the solver's \(S_H(t)\).

  \item \textbf{Entropy match is not a flux match.}
  The 2024 paper notes that \(S_H\) ignores phase.
  A rollout can sit on the right \(S_H\) while \(\Pi(k)\) at small \(k\) has the wrong sign.
  Both are reported. The claim of a learned inverse cascade requires \(\Pi(k)<0\) at small \(k\).
\end{enumerate}

\section{Solver}

Two datasets, both pseudospectral, periodic box \((2\pi)^2\), 2/3 dealiasing, vorticity form
\[
  \partial_t\omega + \mathbf{u}\cdot\nabla\omega = \nu\nabla^2\omega + f - \alpha\omega,
  \qquad
  \nabla^2\psi = -\omega,
  \qquad
  \mathbf{u} = \nabla^\perp\psi.
\]

\begin{center}
\small
\begin{tabular}{lll}
\toprule
& Euler track & Forced track \\
\midrule
Equations & \(\nu = f = \alpha = 0\) & \(\nu > 0\), force at \(k_f = 4\) \\
Grids & \(64^2\), then \(128^2\) & same \\
Time scheme & symplectic, or RK4 & RK4 with integrating factor \\
Initial data & analogues of Runs A, B, C & a spun-up turbulent state \\
Saved field & \(\omega(\mathbf{x},t)\) & \(\omega(\mathbf{x},t)\) \\
\bottomrule
\end{tabular}
\end{center}

The 2022 runs used \(512^2\).
The entropy they measured is carried by a handful of low modes, so \(64^2\) and \(128^2\) are enough to see \(S_H\) fall, provided the dealiasing and the conservation of energy are checked the way the paper checks them: relative drift of \(E\) and of enstrophy below about \(10^{-4}\) on an Euler run.
If the coarse Euler run fails to condense, the grid goes to \(128^2\) before any network is trained.
Forcing uses a single shell so the injection scale is obvious in \(\Pi(k)\).

Each saved frame stores \(\omega\).
Velocity, \(E(\mathbf{k})\), \(S_H\), \(T(k)\), and \(\Pi(k)\) are computed from it in the analysis code, with the same definitions as the papers.

\section{Network}

The model maps one vorticity frame to the next frame at a fixed \(\Delta t\) of a few solver steps, so the target is not contaminated by a single-step truncation error.

Architecture, in order of preference for the first working version:
\begin{enumerate}
  \item A small Fourier neural operator, or a convolutional net of a few layers, with a hard divergence-free projection: the network outputs \(\omega\), and \(\mathbf{u}\) is recovered from \(\psi\).
  \item Parameter count small enough to train on one GPU from the \(64^2\) data in a few hours.
\end{enumerate}

The loss used for the main claim is the pixel loss plus a spectral term,
\[
  \mathcal{L}
  = \|\omega_{\mathrm{pred}}-\omega_{\mathrm{true}}\|_2^2
  + \lambda_E\sum_k\left(\log E_{\mathrm{pred}}(k)-\log E_{\mathrm{true}}(k)\right)^2.
\]
A second training adds the flux term from the 2022 diagnostic,
\[
  \mathcal{L}_\Pi = \sum_{k\le k_f}\big(\Pi_{\mathrm{pred}}(k)-\Pi_{\mathrm{true}}(k)\big)^2,
\]
and the comparison of the two trainings is Prediction 4.
\(\Pi\) is computed from the predicted velocity by the same transfer routine used on the solver output.

No entropy of the weights is in the loss or in the tables.

\section{Diagnostics, matched to the papers}

On every held-out rollout, solver and emulator side by side:

\begin{enumerate}
  \item \(S_H(t)\) from Eq.~\eqref{eq:SH}, modal sum, together with \(S_H/S_{\max}\) and \(S_{\max}=\log_2 M\).
  \item Shell \(S_H(t)\) from the 2024 shell formula, reported for the forced runs.
  \item The fraction of energy in the first few modes, so a Run-B-like state is visible as ``two modes, \(S_H\approx 1\)'' rather than only as a curve.
  \item \(E(k)\) and \(\Pi(k)\) averaged over the late window, as in Fig.~4 of the 2022 paper.
  \item Relative drift of total energy and total enstrophy on Euler rollouts.
  \item One vorticity snapshot of the late state, next to the solver snapshot. \(S_H\) does not show the phase.
\end{enumerate}

Reference lines drawn on the entropy plot: \(0\), the late solver value, and \(\log_2 M\).
A learned Euler run that falls from \(\log_2 M\) toward the solver plateau is the 2022 result.
A learned forced run that stays flat while \(\log_2 M\) moves up under grid refinement is the 2024 result.

\section{Work plan}

\paragraph{Stage 0. Solver and the paper curves.}
Implement the \(64^2\) Euler solver. Reproduce the qualitative 2022 trajectories: \(S_H(t)\) falls for the three initial conditions, the shear-like case ends lowest, and \(\Pi(k)<0\) at small \(k\) in the late window.
Then the forced case: \(S_H(t)\) levels off.
Write down the measured plateaus and \(S_{\max}\).
This stage is finished when those curves exist and energy is conserved on the Euler runs.
No network yet.

\paragraph{Stage 1. One-step model.}
Train on frames from one Euler family and, separately, on the forced series.
Report one-step vorticity error and one-step \(|S_H^{\mathrm{pred}}-S_H^{\mathrm{true}}|\).
This only checks that the map is locally accurate.

\paragraph{Stage 2. The actual comparison.}
Autoregressive rollout over the same horizon as the solver curves from Stage 0.
Overlay \(S_H(t)\).
This is the figure the project is for.
Record which of the three outcomes occurred: condensate ( \(S_H\) tracks the solver, \(\Pi<0\) at low \(k\)), collapse (\(S_H\to 0\)), or thermalization (\(S_H\to\log_2 M\)).

\paragraph{Stage 3. The two predictions that use the papers' structure.}
Train with and without \(\mathcal{L}_\Pi\).
Train on one initial-condition family and roll out the other two.
Repeat the forced run at \(128^2\) and check that the plateau does not scale like \(\log_2 M\).

\paragraph{Stage 4. Write-up.}
One entropy figure per regime, one flux figure, the table of asymptotic \(S_H\) against \(\log_2 M\), and a statement of which predictions held.

\section{What counts as a result}

The project succeeds if Stage 2 produces a rollout whose \(S_H(t)\) follows the solver into an initial-condition-dependent plateau below \(\log_2 M\), with negative low-\(k\) flux, on at least the Euler shear case.
It also succeeds, as a negative result the papers make interpretable, if the rollout entropy goes to \(0\) or to \(\log_2 M\): that says the learned step does not carry the inverse cascade, and Stage 3 then tests whether the flux loss repairs it.

The project fails only if the solver itself does not produce a falling \(S_H\) on the Euler runs.
In that case the grid or the time scheme is wrong, and the network is not the object under test.

\section{Boundaries}

GraphCast and FourCastNet stay out of this plan.
Their spectra are a later application of a diagnostic that has first been checked on a flow where the 2022 and 2024 papers say what \(S_H(t)\) must do.
The three-dimensional Euler thermalization in the 2022 appendix is a useful control if time remains: the same network class, trained on 3D Euler, should raise \(S_H\) rather than lower it.
That control is optional and comes after Stage 2.

\end{document}
